% Generated draft: requires mathematical and source review.
\documentclass[11pt]{article}
\usepackage[margin=25mm]{geometry}
\usepackage{fontspec}
\setmainfont{Latin Modern Roman}
\usepackage{amsmath,amssymb,mathtools,unicode-math}
\setmathfont{Latin Modern Math}
\usepackage{enumitem,needspace,xurl,hyperref}
\setlength{\parindent}{0pt}
\setlength{\parskip}{5pt}
\setlength{\emergencystretch}{3em}
\newcommand{\N}{\mathbb{N}}
\newcommand{\Z}{\mathbb{Z}}
\newcommand{\Q}{\mathbb{Q}}
\newcommand{\R}{\mathbb{R}}
\newcommand{\C}{\mathbb{C}}
\newcommand{\F}{\mathbb{F}}
\newcommand{\A}{\mathbb{A}}
\newcommand{\PP}{\mathbb{P}}
\newcommand{\E}{\mathbb{E}}
\newcommand{\eps}{\varepsilon}
\newcommand{\dd}{\,\mathrm{d}}
\begin{document}
\section*{UnsolvedMath ID: 30004823}
\textbf{Source problem code:} catalogue entry
\subsection{Definitions and mathematical statement}
Let \(S_n\) denote the symmetric group on \(n\) letters. For a partition \(\mu\vdash n\), let \(S^{\mu}_{\Q}\) denote the corresponding rational Specht module. For a finite-dimensional rational \(S_n\)-representation \(W\), define its Specht multiplicity by
\[
[W:S^{\mu}_{\Q}]
:=\dim_{\Q}\operatorname{Hom}_{\Q[S_n]}(S^{\mu}_{\Q},W).
\]
The length \(\ell(\lambda)\) of a partition \(\lambda\) is its number of positive parts.

The supplied record describes an ideal \(I\) in an inverse-limit ring \(\Lambda\) of symmetric polynomials, a sequence of associated symmetric varieties \(V_n\), and a partition \(\{\lambda\}_n\) indexing an \(S_n\)-representation for sufficiently large \(n\). It does not specify the coefficient ring, the inverse-limit convention and transition maps, the construction or ambient field of \(V_n\), or the rule defining \(\{\lambda\}_n\). These are unresolved parts of the formulation. In particular, no partition-padding convention or real-versus-complex interpretation of the varieties is imposed here. The cohomology theory and the hypotheses ensuring finite-dimensional rational cohomology with the required \(S_n\)-action also need confirmation.

Subject to those missing definitions and hypotheses, the intended conjecture is the following: for every ideal \(I\) allowed by the intended construction, every fixed partition \(\lambda\), and every fixed integer \(p\geq 0\), there exist an integer \(N\) and a polynomial \(P_{I,\lambda,p}(t)\) such that, for every integer \(n\geq N\), the partition \(\{\lambda\}_n\) is defined and
\[
m_{I,\lambda,p}(n)
:=\bigl[H^p(V_n;\Q):S^{\{\lambda\}_n}_{\Q}\bigr]
=P_{I,\lambda,p}(n).
\]
The threshold and polynomial may depend on \(I\), \(\lambda\), and \(p\). This is a conditional statement of the supplied mathematical target; the record does not determine a complete, unambiguous class of instances.

The supplied excerpt also reports a special case, described as joint work with D. Perrucci. It concerns a principal ideal \(I=(f)\), where \(f\) is a finite linear combination of elementary symmetric functions. The extracted sum suggests
\[
f=\sum_{i=0}^{d}a_i\sigma_i\in\Lambda,
\]
where \(\sigma_i\) denotes the elementary symmetric function of degree \(i\), with the usual convention \(\sigma_0=1\). However, the accompanying extracted range is literally \(0\leq i\geq d\); the finite range \(0\leq i\leq d\) is therefore a plausible reading requiring source confirmation. The coefficient domain and any further hypotheses on \(f\) are not supplied. In the intended setting, the excerpt reports that for every fixed partition \(\lambda\), for all sufficiently large \(n\),
\[
m_{(f),\lambda,0}(n)=
\begin{cases}
0,&\ell(\lambda)>1,\\
c_{f,\lambda},&\ell(\lambda)=1,
\end{cases}
\qquad c_{f,\lambda}\in\Z_{\geq0}.
\]
It presents this as verifying the conjecture for these principal ideals when \(p=0\). This is a supplied special-case claim, not an independently verified theorem here. If the empty partition is allowed, the displayed cases do not specify its treatment.
\subsection{Short English statement}
For each fixed ideal, partition, and cohomological degree, the question asks whether the corresponding Specht multiplicity eventually equals a polynomial in the number of variables. The supplied record reports a special case in degree zero, but essential definitions of the varieties and partition indexing remain unspecified. This statement-only draft requires human mathematical review.
\subsection{Sources}
\begin{itemize}
\item Supplied UnsolvedMath local registry record, permanent ID \(30004823\), title ``Polynomial Specht Multiplicities in Symmetric-Variety Cohomology,'' exact source problem code \texttt{catalogue entry}; collection: Oberwolfach Reports Open Problems. No proposer is supplied.
\item Supplied report citation: \emph{MFO-RIMS Tandem Workshop: Symmetries on Polynomial Ideals and Varieties (2022)}, Oberwolfach Reports, pp. 2365--2396, DOI \texttt{10.4171/owr/2021/43}. \url{https://doi.org/10.4171/owr/2021/43}. The registry gives proposed year 2021, whereas the supplied citation title contains 2022; this discrepancy is unresolved.
\item Supplied literature source labeled ``Basu--Riener'': \url{https://academic.oup.com/imrn/article/2020/7/2054/4989853}. No fuller bibliographic identification or precise supporting result is provided in the record.
\item The extracted report passage labels the conjecture ``Conjecture 1. [2]'' and the special case ``Theorem 1. [1].'' The bibliography entries corresponding to [1] and [2] are not supplied. The special case is described as joint work with D. Perrucci; the other participant is not identified in this record.
\item Literature assessment date: \texttt{2026-08-22}. The supplied dated curation labels the problem ``partially solved,'' reports polynomial bounds and special cases including the stated \(p=0\) family, and says that the general eventual-polynomial multiplicity conjecture remains open. Its listed sources are the OWR report and Basu--Riener links above. This assessment is supplied dated curation derived from the OWR record, not independently verified here; its time-sensitive status and cited claims require independent checking.
\item The registry labels its statement \texttt{corrected\_verified} and says that the official report was checked directly and extraction artifacts were removed. This is a supplied verification claim, not a source visit or verification performed for this draft.
\end{itemize}
\end{document}
